\section{Reuse the exact eight-tap calculation}
Let $x_c[n]$ be input frame $n$ in channel $c$. For this file operation define
$x_c[n]=0$ when $n<0$. The coefficient codes $h_0$ through $h_7$ are ordered from
oldest to newest input, as in the earlier FIR interface. For each output frame,
\[
  S_c[n]=\sum_{j=0}^{7}h_j\,x_c[n-7+j],\qquad
  R_c[n]=\left\lfloor\frac{S_c[n]+16384}{32768}\right\rfloor.
\]
Store $R_c[n]$ clamped to the signed sixteen-bit range. Count a clip only if the
rounded result is strictly outside that range. Sum exact signed products before
rounding once; neither per-product rounding nor intermediate saturation is equivalent.
The eight-term sum can require more than thirty-two signed bits.

Produce exactly one output frame per input frame. The seven zero history values set
the beginning condition independently for each channel. Do not append a seven-frame
decay after the last input: this interface deliberately keeps the original frame count
and duration. A longer convolution with an emitted tail is a different operation.

The wrapper may reuse \texttt{fir\_q15} in blocks or implement another exact method.
The earlier kernel returns its processed count in \texttt{v0}; this wrapper returns a
status there. Retain the block length and accumulate the earlier kernel's clip counts
without confusing either result with the wrapper's complete frame count.
